% --- Back cover ---
% Large PERSONIX wordmark, no logo image. Cream paper inherited globally.
% Absolute (tikz overlay) positioning so the footer sticks to the bottom.
\clearpage
\thispagestyle{empty}
\begin{tikzpicture}[remember picture, overlay]
  % --- Wordmark ---
  \node[anchor=north, font=\sffamily\bfseries\fontsize{48}{52}\selectfont]
    at ([yshift=-26mm]current page.north) {PERSONIX};
  \node[anchor=north, font=\rmfamily\itshape\large]
    at ([yshift=-44mm]current page.north) {Uzlaşmasız Dəyişiklik};
  \node[anchor=north, font=\rmfamily\small,
        text width=\paperwidth-50mm, align=center]
    at ([yshift=-54mm]current page.north)
    {Senzura Edilə Bilməyən və Korrupsiyaya Uğramayan Mərkəzsizləşdirilmiş Reputasiya Şəbəkəsi};
  % --- Body ---
  \node[anchor=north, text width=\paperwidth-50mm, align=justify, font=\small]
    at ([yshift=-72mm]current page.north)
    {%
      Ünsiyyət yavaş, qərarlar yerli olduqda və avtoritet idarə etdiyi
      şəhərdə yaşadıqda dövlət işləyirdi. Bugünkü şəbəkələr hər üç
      fərziyyəni tərsinə çevirir --- və onların üzərində qurulmuş institutlar
      artıq idarə etdikləri dünyaya uyğun gəlmir. Bu kitab uyğun gələn bir
      aləti təsvir edir: mərkəzsizləşdirilmiş reputasiya şəbəkəsi, harada ki,
      kimlik saxlamaq üçün sənindir, həqiqət açıq şəkildə audit oluna bilər
      və rüşvət reputasiya intiharıdır.

      \smallskip
      Manifest deyil. Proqnoz deyil. Dövlətin varisinin necə qurula
      biləcəyinə dair işlək bir plan --- könüllü şəkildə, hər dəfə bir
      bəyanatla.
    };
  % --- Footer (anchored to the bottom of the page) ---
  \node[anchor=south, font=\sffamily\footnotesize,
        text width=\paperwidth-50mm, align=center]
    at ([yshift=12mm]current page.south)
    {%
      Pavel Kudrna\quad\textbullet\quad Praqa, 2026\par
      \href{https://personix.org}{personix.org}\par
      Creative Commons Attribution-ShareAlike 4.0 International
      (CC BY-SA 4.0) lisenziyası altında.
    };
\end{tikzpicture}%
\mbox{}%
\clearpage
